\section{总结与延伸}

本讲用一个极简问题重构了 offline RL：若完整轨迹本来就是序列，为什么不直接学习“在目标回报与历史条件下，下一个动作是什么”？这个问题带来了一套清晰方法，也带来一套必须诚实面对的假设。

\begin{enumerate}
  \item \textbf{问题表述是主要创新}：将 RL 写成 conditional sequence modeling，而非仅替换 policy backbone；
  \item \textbf{return-to-go 是控制 token}：训练时由完整日志回算，推理时作为 desired behavior 条件；
  \item \textbf{监督目标改善可训练性}：MSE/CE 避免显式 Bellman backup，但不自动进行 policy improvement；
  \item \textbf{history 是有意设计}：context 可处理 partial observability 和长期事件，也可能记住数据偏差；
  \item \textbf{rollout 才是最终考试}：requested return、训练 loss 与 realized environmental return 必须分开；
  \item \textbf{强简单 baseline 必不可少}：percent BC 揭示数据筛选与条件建模各自贡献；
  \item \textbf{sparse reward 不等于 exploration}：离线日志可传播未来成功信息，但在线 cold start 仍难；
  \item \textbf{critic 实验展示接口弹性}：同一 causal model 可预测 action 或 return，却不自动构成完整 actor-critic；
  \item \textbf{外推是观察，不是保证}：极端 target 会饱和，偶发超越 logged best 需要重复验证；
  \item \textbf{未来更可能是 hybrid}：sequence model 与 pessimism、uncertainty、online data collection 可以组合，但需重新做因果与稳定性审计。
\end{enumerate}

\begin{importantbox}{课程作业：设计一组能区分“条件建模”与“数据筛选”的实验}
选择一个 offline-RL benchmark，固定相同 encoder、parameter count 和训练预算，比较：普通 BC、percent BC、Decision Transformer、$K=1$ DT 与带 conservative critic 的 DT。对每个模型报告数据量曲线、target-return sweep、realized return、OOD action rate、五个随机种子及失败轨迹。最后回答：提升来自 return condition、history、更多有效数据，还是 critic filtering？
\end{importantbox}

\begin{warningbox}{最终自检：五个不能混为一谈的结论}
\textbf{训练稳定}不等于\textbf{闭环稳定}；\textbf{条件可控}不等于\textbf{目标可达}；\textbf{长期相关}不等于\textbf{因果 credit}；\textbf{离线成功}不等于\textbf{在线探索成功}；\textbf{统一 token 接口}不等于\textbf{所有任务共享同一难度}。
\end{warningbox}

\subsection{拓展阅读}

这一组材料从原始方法、轨迹建模、保守离线 RL、数据集到 probabilistic inference 逐层展开，适合按下列顺序阅读：

{\small
\begin{itemize}[nosep,leftmargin=*]
  \item \href{https://arxiv.org/abs/2106.01345}{\textbf{Decision Transformer}}：方法定义、benchmark、ablation 与 Key-to-Door 实验。
  \item \href{https://sites.google.com/berkeley.edu/decision-transformer}{\textbf{Project Page}} 与 \href{https://github.com/kzl/decision-transformer}{\textbf{Official Code}}：可视化、数据处理与 rollout 实现。
  \item \href{https://arxiv.org/abs/2106.02039}{\textbf{Trajectory Transformer}}：将 dynamics、reward 与 action 更完整地写成轨迹生成问题。
  \item \href{https://arxiv.org/abs/2006.04779}{\textbf{Conservative Q-Learning}}：offline RL 中 pessimism 与 distributional support 的另一条路线。
  \item \href{https://arxiv.org/abs/2004.07219}{\textbf{D4RL}}：offline-RL dataset 与 normalized evaluation protocol。
  \item \href{https://arxiv.org/abs/1805.00909}{\textbf{RL and Control as Probabilistic Inference}}：optimality variable、maximum entropy 与 inference 视角。
\end{itemize}
}

\end{document}
